%% sec_introduction.tex
%% ─────────────────────────────────────────────────────────────────────────────
%% Section 1 — Introduction
%% ─────────────────────────────────────────────────────────────────────────────

\section{Introduction}\label{sec:intro}

Due to the intrinsically probabilistic nature of fundamental particle
interactions and the composite structure of the proton, accurately modelling
the dynamics of proton collisions at the Large Hadron Collider (LHC)
constitutes one of the major challenges in high-energy physics
(HEP)~\cite{HSFPhysicsEventGeneratorWG:2020gxw}.  Physics event generators
encode the necessary theoretical tools in computational form, enabling
quantitative predictions for a given model and process via Monte Carlo (MC)
techniques.

For a given process, the generator samples a large number of random
phase-space points (PSPs) and evaluates the squared matrix element \ME at each
point.  As the LHC moves into its High-Luminosity phase (\hilumi), the
simulation demand grows super-linearly while the computing budget remains
flat---a gap that cannot be closed by conventional CPU improvements
alone~\cite{HSFPhysicsEventGeneratorWG:2020gxw}.  The matrix element (ME)
evaluation step is the dominant computational bottleneck, accounting for
30--40\% of total generator CPU time in multi-jet processes~\cite{Valassi:2021ljk}.

This bottleneck has driven substantial parallelism efforts.  The \amcnlo
(\mgfive) CUDACPP plugin~\cite{Valassi:2025xfn} exploits GPU and
vectorised-CPU architectures to deliver speedups of
$100\times$--$1000\times$ over single-threaded Fortran on complex processes.
Field-Programmable Gate Arrays (FPGAs) offer an orthogonal route: custom
fixed-function datapaths designed around the exact computation required, with
typically far lower power consumption than GPUs.  Barbone et
al.~\cite{Carrazza:2021gpx} demonstrated a ${\sim}100\times$ CPU speedup for
the $e^+e^-\to\mu^+\mu^-$ matrix element on a Xilinx Alveo card.  The AMD
Versal Adaptive Compute Acceleration Platform (ACAP) extends this further by
co-integrating a programmable logic (PL) fabric with an array of specialised
\aie cores~\cite{Barbone:2023mul}.

In this paper we describe the first full-array implementation of a
leading-order ME pipeline on the Versal VCK190, targeting the
\ggttg process generated by \mgfive.  Our central contribution is an
\textit{architectural} one: we show how the 16\KB program-memory (PM) limit of
a single \aie tile fundamentally constrains the design, and we present a
systematic methodology---driven by per-function PM profiling, a formal cascade
contract, and a novel wavefunction-token protocol---that maps the complete
16-diagram \ggttg computation onto 80 parallel cascade pipelines occupying all
400 \aie tiles of the device.

The specific contributions of this work are:
\begin{enumerate}
  \item \textbf{Wavefunction-token cascade pipeline architecture.}
    A five-stage pipeline (\texttt{K1}$\to$\texttt{K2a}$\to$\texttt{K2b}%
    $\to$\texttt{K3}$\to$\texttt{K4}) that decomposes 16 Feynman diagrams
    across five \aie tiles.  Tiles communicate via a 384-bit cascade token
    carrying external wavefunctions, precomputed VVV propagators, and partial
    JAMP accumulators.

  \item \textbf{Program-memory-driven partitioning methodology.}
    We profile every HELAS function's PM footprint, derive the \textit{VVV
    Exclusion Principle} (\texttt{VVV1P0\_1} and \texttt{VVV1\_0} cannot
    coexist in one tile), and apply a lazy-evaluation optimisation that moves
    \texttt{FFV1P0\_3} from \texttt{K1} to \texttt{K3}, eliminating the last
    PM overflow.

  \item \textbf{Deterministic cascade contract.}
    Formal rules---all kernels share the same \texttt{BATCH}$\times$32 loop
    skeleton; cascade writes are unconditional; token counts are statically
    matched---that guarantee deadlock-free operation without flow-control
    hardware.

  \item \textbf{HELAS-to-AIE adaptation.}
    Porting of the full MadGraph HELAS library to \texttt{aie::vector<cfloat,8>}
    intrinsics, including complex-division elimination, vectorised spinor cores,
    and binary-indexed helicity caching.

  \item \textbf{Scalable packet-switched array and full PS--PL--AIE system.}
    Scale-up to 80 independent pipelines via
    $10\times\texttt{pktsplit}\langle 8\rangle$, embedded in a three-layer
    system: ARM A72 PS (RAMBO phase-space generation + XRT), PL (10$\times$ HLS
    MM2S/S2MM data movers), and the \aie array.

  \item \textbf{Physics equivalence validation.}
    Diagram-by-diagram JAMP comparison against the \mgfive double-precision
    reference, documenting the float32 tolerance.
\end{enumerate}

The remainder of the paper is structured as follows.
Section~\ref{sec:background} reviews \mgfive and the \ggttg process.
Section~\ref{sec:platform} describes the Versal VCK190 target hardware.
Section~\ref{sec:architecture} presents the cascade pipeline architecture.
Section~\ref{sec:implementation} details the HELAS adaptation and system
integration.
Section~\ref{sec:results} reports program-memory, precision, throughput, and
energy results.
Section~\ref{sec:conclusion} concludes and discusses future directions.
